\documentclass[10pt]{article}
\usepackage[T1]{fontenc}
\usepackage{lmodern,amsmath,amssymb,amsthm,microtype}
\usepackage[margin=0.85in]{geometry}
\usepackage[colorlinks=true,urlcolor=blue]{hyperref}
\hypersetup{
  pdftitle={An All-Index Arithmetic Minimum Formula for OEIS A069862},
  pdfauthor={Qichao Wang; Daoyu Dong},
  pdfsubject={Universal existence and exact least concatenation lengths}
}
\newtheorem{theorem}{Theorem}
\newtheorem{lemma}{Lemma}
\newtheorem{proposition}{Proposition}
\newtheorem{corollary}{Corollary}
\newcommand{\A}{\mathrm{A069862}}
\title{An All-Index Arithmetic Minimum Formula for OEIS A069862}
\author{Qichao Wang\\{\small Hebei University of Technology, Tianjin, China}\and
Daoyu Dong\\{\small University of Electronic Science and Technology of China, Chengdu, China}}
\date{October 22, 2026}
\begin{document}
\maketitle
\begin{abstract}
For OEIS A069862, let $a(n)$ be the least positive $k$ for which the decimal concatenation of $n+1,\ldots,n+k$ is divisible by $n$. We give an exact finite arithmetic formula for this minimum for every positive integer $n$. In each digit interval, multiplicative-order phases reduce the coprime part to quadratic congruences, while the 2- and 5-parts reduce to a linear endpoint condition after at most three exceptional terms. Modular square roots and CRT give all admissible progressions, and ceilings select their first endpoints. Constructive prime-power lifting proves existence, infinitely many witnesses, and the bound $a(n)<10^n-n$ for $n\ge2$, making the formula's range explicit. The formula retains a finite minimum over arithmetic candidates; it is not an elementary expression without minimization. We also retain two infinite families of simpler exact identities. The known classification of the indices satisfying $a(n)=2$ was recorded by David Wasserman (2005) in OEIS A069860 and is excluded from our novelty claims.
\end{abstract}
\section{Definition and a modular recurrence}
Let $\ell(t)$ denote the number of base-10 digits of $t\ge1$. Define
\[
 C_n(0)=0,\qquad C_n(k)=10^{\ell(n+k)}C_n(k-1)+(n+k)\quad(k\ge1).
\]
Whenever $C_n(k)\equiv0\pmod n$, the sequence value $a(n)$ is the least such $k$. For $n>1$, $C_n(1)\equiv1\pmod n$, and $a(1)=1$.
\begin{proposition}[Constant-width reduction]\label{prop:width}
Put $D=\ell(n+1)$ and $L=10^D-n-1$. For $1\le k\le L$, write $R_k=C_n(k)\bmod n$. Then all the concatenated blocks have $D$ digits. Consequently,
\[
 \begin{array}{rcll}
 R_k&\equiv& k &(10^D\equiv0\pmod n),\\
 R_k&\equiv&k(k+1)/2 &(10^D\equiv1\pmod n),\\
 R_k&\equiv&\lceil k/2\rceil &(10^D\equiv-1\pmod n).
 \end{array}
\]
The triangular expression is an integer before reduction, even for even $n$.
\end{proposition}
\begin{proof}
In this range $n+i<10^D$ and $n+i\ge n+1\ge10^{D-1}$, so the recurrence modulo $n$ is $R_k\equiv B R_{k-1}+k$ with $B=10^D$ and $R_0=0$. When $B\equiv0$, the result is $k$; when $B\equiv1$, it is $\sum_{i=1}^k i=k(k+1)/2$. When $B\equiv-1$, induction gives $R_{2t}\equiv t$ and $R_{2t+1}\equiv t+1$.
\end{proof}
\section{Existence for every index}\label{sec:existence}
The local lifting method used for Champernowne prefixes in the repository's
A075002 manuscript~\cite{prefix} also applies to a prefix starting at $n+1$.
The essential extra observation is that when $3\mid n$, the removed
triangular sum $n(n+1)/2$ vanishes modulo~3. We give the full argument
so that no certificate from another sequence is needed.

\begin{theorem}[All-index existence]\label{thm:existence}
For every positive integer $n$, there are infinitely many $t>n$ such that
\[
 n\mid F_n(t),\qquad F_n(t)=(n+1)\Vert(n+2)\Vert\cdots\Vert t.
\]
In particular, $a(n)$ exists for every $n\ge1$.
\end{theorem}

First put $F_n(n)=0$. In a fixed $d$-digit interval, let $B=10^d$ and
suppose $t+1,\ldots,t+L$ all lie in that interval. Appending these blocks
gives, for $L\ge0$,
\begin{equation}\label{eq:shift}
 F_n(t+L)-F_n(t)=(B^L-1)F_n(t)+tS_L+T_L,
\end{equation}
where $S_L=\sum_{j=0}^{L-1}B^j$ and
$T_L=\sum_{j=1}^L jB^{L-j}$. The telescoping identity
\begin{equation}\label{eq:weighted}
 (B-1)T_L=BS_L-L
\end{equation}
will be used only when the indicated division is legitimate.

\begin{lemma}[Two local shifts]\label{lem:shifts}
Assume the fixed-width condition in \eqref{eq:shift}.
\begin{enumerate}
\item Let $p\ge7$ be prime with $p\nmid B(B-1)$, and let
$h=\operatorname{ord}_p(B)$. If $L=chp^a$, $a\ge0$, and $p\nmid c$, then
\[
 F_n(t+L)-F_n(t)\equiv-\frac{chp^a}{B-1}\pmod{p^{a+1}}.
\]
For $e\ge1$, every multiple of
$p^e\operatorname{ord}_{p^e}(B)$ preserves both $F_n(t)$ and $t$
modulo $p^e$, provided the appended blocks stay in the interval.
\item If $a=v_3(L)\ge1$, then
\[
 F_n(t+L)-F_n(t)\equiv Lt+\frac{L(L+1)}2\pmod{3^{a+1}}.
\]
Every multiple of $2\cdot3^e$ preserves both $F_n(t)$ and $t$
modulo $3^e$, under the same interval condition.
\end{enumerate}
\end{lemma}
\begin{proof}
For the first assertion, the binomial theorem (or LTE) gives
$B^{chp^a}\equiv1\pmod{p^{a+1}}$, since $B^h\equiv1\pmod p$.
Thus $S_L\equiv0$ because $B-1$ is a unit. Equations
\eqref{eq:shift}--\eqref{eq:weighted} give the formula.
For a multiple of $p^e\operatorname{ord}_{p^e}(B)$, the same identities
give $B^L=1$, $S_L=T_L=0$ modulo $p^e$ and $p^e\mid L$.

For the second assertion write $B=1+c$ with $9\mid c$. Expansion gives
\[
 S_L-L=\sum_{j\ge1}c^j\binom{L}{j+1},\qquad
 T_L-\frac{L(L+1)}2=\sum_{j\ge1}c^j\binom{L+1}{j+2}.
\]
The valuations of the summands are respectively at least
$2j+a-v_3(j+1)$ and $2j+a-v_3((j+1)(j+2))$, both at least $a+1$.
Here $v_3(j+1)\le j$, and
$v_3((j+1)(j+2))\le2j-1$ for $j\ge1$ (check $j=1$ directly;
for $j\ge2$, $3^{2j}>(j+1)(j+2)$).
Also $B^L\equiv1\pmod{3^{a+2}}$ by LTE. Substitute in
\eqref{eq:shift}. If $2\cdot3^e\mid L$, the displayed shift is
zero modulo $3^e$, and $t$ is preserved as well.
\end{proof}

\begin{proof}[Proof of Theorem~\ref{thm:existence}]
The case $n=1$ is immediate. Write
\[
 n=R\,3^e\prod_{i=1}^s p_i^{e_i},\qquad
 R=2^\alpha5^\beta,\qquad 7\le p_1<\cdots<p_s.
\]
Let $D=\operatorname{lcm}_i\operatorname{ord}_{p_i^{e_i}}(10)$,
with $D=1$ for an empty product. Choose $d\equiv1\pmod D$ sufficiently
large, so $d\ge\max(\alpha,\beta)$ and $10^{d-1}>n$.
Then $B=10^d$ is zero modulo $R$ and is congruent to 10 modulo
each $p_i^{e_i}$. All step sizes below are consequently independent of $d$.

Define $H_0=2\cdot3^e$ if $e>0$, and $H_0=1$ otherwise; recursively put
\[
 H_i=\operatorname{lcm}\bigl(H_{i-1},
       p_i^{e_i}\operatorname{ord}_{p_i^{e_i}}(10)\bigr).
\]
For $1\le a<e$ let $L_a=\operatorname{lcm}(2\cdot3^a,R)$.
For $0\le a<e_i$ put
\[
 J_{i,a}=\operatorname{lcm}
   \bigl(\operatorname{ord}_{p_i}(10)p_i^a,H_{i-1},R\bigr),
 \qquad M_0=R3^{\min(e,1)}.
\]
An explicit displacement bound is
\begin{equation}\label{eq:capacity}
 W=2\sum_{a=1}^{e-1}L_a+
       \sum_{i=1}^s(p_i-1)\sum_{a=0}^{e_i-1}J_{i,a},
 \qquad 9\cdot10^{d-1}>M_0+W.
\end{equation}
Empty sums are zero. Increase $d$ within its progression to ensure
\eqref{eq:capacity}, and let $t_0$ be the least multiple of $M_0$
at least $10^{d-1}$. Every candidate below lies between $t_0$ and
$t_0+W<10^d$, which proves every required fixed-width condition.

The last block gives $F_n(t_0)\equiv t_0=0\pmod R$.
If $e>0$, then $3\mid n,t_0$, so
\[
 F_n(t_0)\equiv\sum_{j=n+1}^{t_0}j
   =\frac{t_0(t_0+1)-n(n+1)}2\equiv0\pmod3.
\]
At a root modulo $3^a$, $1\le a<e$, choose among endpoints
$t+jL_a$, $j=0,1,2$. Write $L_a=2c3^a$, where $3\nmid c$.
Lemma~\ref{lem:shifts} gives
\[
 \frac{F_n(t+jL_a)-F_n(t)}{3^a}
       \equiv jc(2t+1)\pmod3.
\]
Since $3\mid F_n(t)$ and $3\mid n$, the triangular identity implies
$t\equiv0$ or $2\pmod3$. The slope $c(2t+1)$ is therefore a unit.
Exactly one choice raises the root to $3^{a+1}$. All endpoints remain
multiples of $R$, preserving the last-block condition. This installs
$R3^e$; the period $H_0$ preserves the completed 3-part.

Now process $p_i$ in increasing order. Every prime divisor of
$p_j^{e_j}\operatorname{ord}_{p_j^{e_j}}(10)$ is at most $p_j$,
since the order divides $p_j^{e_j-1}(p_j-1)$. Thus $p_i\nmid RH_{i-1}$.
Consequently $J_{i,a}=chp_i^a$ with $h=\operatorname{ord}_{p_i}(10)$
and $p_i\nmid c$. At a root modulo $p_i^a$ (vacuous for $a=0$),
Lemma~\ref{lem:shifts} shows that the endpoints
$t+jJ_{i,a}$, $0\le j<p_i$, sweep all next $p_i$-adic digits,
with unit normalized slope $-ch/(B-1)$ modulo $p_i$.
Choose the unique root modulo $p_i^{a+1}$. The multiples of $H_{i-1}$
preserve all earlier odd prime powers, and those of $R$ preserve
the last-block condition. After installing $p_i^{e_i}$, the period
$H_i$ preserves it too. Induction yields $n\mid F_n(t)$.

We have $n<t<10^d$, hence the explicit bound $a(n)\le t-n<10^d-n$.
There are arbitrarily large admissible $d$; their disjoint digit
intervals give infinitely many endpoints. This proves the theorem.
\end{proof}

\begin{corollary}[A uniform bound]\label{cor:uniform}
For every $n\ge2$,
\[
 1\le a(n)<10^n-n.
\]
\end{corollary}
\begin{proof}
Retain the notation in the proof, and put $K=\operatorname{lcm}(R,H_s)$.
Every $L_a$ and $J_{i,a}$ divides $K$, and
\[
 K\le 2n\prod_i\varphi(p_i^{e_i})\le2n^2.
\]
The elementary inequalities $p^f\ge1+f(p-1)$ and
$\prod_j(1+x_j)\ge1+\sum_j x_j$ for nonnegative $x_j$ give
\[
 2\max(e-1,0)+\sum_i(p_i-1)e_i\le n-1.
\]
Thus $W\le2n^2(n-1)$ and $M_0\le n$.
Let $b=\lceil\log_2 n\rceil$ and choose
\[
 d=1+D\lceil b/D\rceil.
\]
Then $d\ge b+1$, so $10^{d-1}\ge10^b\ge n^3>n$ and
$9\cdot10^{d-1}>n+2n^2(n-1)$. All conditions in the construction hold.
Also $D\le\prod_i\varphi(p_i^{e_i})\le n-1$ (including the
empty-product case). If $D\ge b$, then $d=D+1\le n$.
If $D<b$, then $d\le b+D\le2b-1\le n$; the last inequality holds
for $n\ge3$, and $n=2$ falls in the first case.
The constructed endpoint therefore satisfies $n<t<10^d\le10^n$.
\end{proof}

\begin{proposition}[Exact minimum for every index]\label{prop:algorithm}
Every value $a(n)$ can be determined in at most $n^3$ modular
single-block transitions, together with at most $n$ binary block jumps.
\end{proposition}
\begin{proof}
The case $n=1$ is immediate. For $n\ge2$, process digit intervals
in increasing order, starting with the integer $n+1$ and stopping before
$10^n$. Within a $d$-digit interval, the state
$(F_n(t)\bmod n,t\bmod n)$ lies in a set of size $n^2$ and evolves by
\[
 (x,y)\longmapsto(10^d x+y+1,\ y+1)\pmod n.
\]
Test the residue after each actual appended block. If it is zero,
the first such endpoint gives the exact minimum. If a state repeats
before a zero occurs, determinism implies that the remaining orbit
in this interval has no zero. Jump over the unused part of the interval
by binary exponentiation of its affine transition and proceed to the
next digit width. A repeated state occurs within $n^2$ transitions,
irrespective of whether $10^d$ is invertible modulo $n$.
There are at most $n$ digit intervals, and Corollary~\ref{cor:uniform}
guarantees a zero before they end. No guessed numerical search limit
or assumed general minimum formula is used.
\end{proof}


\section{An all-index arithmetic formula for the minimum}\label{sec:minimum-formula}
We now determine the minimum itself for every index. The expression below
uses multiplicative orders, modular square roots, CRT, ceilings, and a
finite minimum over explicitly constructed residue classes. It is a
\emph{finite arithmetic formula}, not an elementary expression without
finite minimization. It does not enumerate concatenation lengths, and no
unknown earlier sequence values occur in it.

\subsection{Quantities depending only on the index}
Fix $n\ge2$ and write
\[
 n=Rm,\qquad R=2^\alpha5^\beta,\qquad \gcd(m,10)=1,
 \qquad D=\ell(n+1).
\]
For each $D\le d\le n$, put
\[
 B_d=10^d,\qquad t_d=\max(n,10^{d-1}-1),\qquad
 L_d=10^d-1-t_d.
\]
The actual $d$-digit endpoints are $t_d+1,\ldots,t_d+L_d$.
For any $B>1$ and integer $r\ge0$, define the integers
\begin{equation}\label{eq:ST-closed}
 S_B(r)=\frac{B^r-1}{B-1},\qquad
 T_B(r)=\frac{B^{r+1}-(r+1)B+r}{(B-1)^2}.
\end{equation}
In particular $S_B(0)=T_B(0)=0$. An explicit, nonrecursive expression
for the prefix preceding this digit interval is
\begin{equation}\label{eq:boundary-sum}
 c_d=\sum_{e=D}^{d-1}
 \bigl(t_e S_{B_e}(L_e)+T_{B_e}(L_e)\bigr)
 10^{\sum_{f=e+1}^{d-1} fL_f}.
\end{equation}
The empty sum is zero. This is simply the sum of the complete
fixed-width words, shifted by the number of digits following them.
Thus $c_d=F_n(t_d)$. Computation may replace $c_d$ by its residue modulo $n$.

For the remainder of this subsection fix $d$ and abbreviate
$B=B_d,t=t_d,c=c_d$. Put
\[
 E=\max\left(1,\left\lceil\alpha/d\right\rceil,
                  \left\lceil\beta/d\right\rceil\right),\quad
 h=\operatorname{ord}_m(B),\quad s=S_B(h),\quad v=T_B(h),
\]
where $h=1$ when $m=1$. Equivalently, for $m>1$,
$h=\operatorname{ord}_m(10)/\gcd(\operatorname{ord}_m(10),d)$.
For $0\le j<h$ define
\[
 V_j=B^jc+tS_B(j)+T_B(j),
\]
and the quadratic
\begin{equation}\label{eq:phase-quadratic}
 Q_j(z)=hsz^2+\bigl(2(t+j)s+2v-hs\bigr)z+2V_j.
\end{equation}
Let $\mathcal R_j$ be a set of disjoint residue classes $(r,q)$,
with $q\mid m$, whose union is precisely the roots of
$Q_j(z)\equiv0\pmod m$. The explicit square-root construction below
gives at most $2^{\omega(m)}$ such classes; for $m=1$ use
$\mathcal R_j=\{(0,1)\}$.

Write $[x]_u$ for the least nonnegative residue modulo $u$, with
$[x]_1=0$ and all modular-inverse expressions modulo 1 interpreted as 0.
Define
\[
 \rho=[-B(B-1)^{-1}]_R,\qquad g=\gcd(h,R),\qquad u=R/g.
\]
Discard a phase $j$ if $g\nmid\rho-t-j$. For every other phase put
\[
 b_j=\left[\frac{\rho-t-j}{g}(h/g)^{-1}\right]_u.
\]
For each $(r,q)\in\mathcal R_j$, the CRT combination and the first
eligible endpoint are the explicit expressions
\begin{align}
 w&=r+q[(b_j-r)q^{-1}]_u, & P&=hqu,\label{eq:crt-candidate}\\
 x_0&=t+j+hw, &
 X_{d,j,r,q}&=x_0+P\left\lceil\frac{t+E-x_0}{P}\right\rceil.
 \label{eq:first-candidate}
\end{align}
All inverses here exist in the stated moduli.
Define the finite sets
\begin{align*}
 \mathcal E_d={}&\{t+r:1\le r<E,\ r\le L_d,
          \ n\mid B^rc+tS_B(r)+T_B(r)\},\\
 \mathcal C_d={}&\{X_{d,j,r,q}:0\le j<h,\ (r,q)\in\mathcal R_j,
       \ g\mid\rho-t-j,\ X_{d,j,r,q}\le10^d-1\}.
\end{align*}

\begin{theorem}[All-index minimum formula]\label{thm:minimum-formula}
For $n\ge2$, the following finite set is nonempty, and
\begin{equation}\label{eq:minimum-formula}
 \boxed{\displaystyle
 a(n)=\min\bigcup_{d=D}^{n}(\mathcal E_d\cup\mathcal C_d)-n.}
\end{equation}
Together with $a(1)=1$, this determines the entire sequence. For each
digit width there are at most three exceptional endpoints and at most
$h\,2^{\omega(m)}$ progression candidates. All quantities in the formula
are specified by $n$, with no test of successive lengths up to $a(n)$.
\end{theorem}

\subsection{Proof of completeness and minimality}
For a fixed width, write $G(r)=B^rc+tS_B(r)+T_B(r)$, initially for
nonnegative integers $r$. This is the fixed-width continuation of the
prefix; it equals the actual concatenation whenever $r\le L_d$.
Multiplying by $(B-1)^2$ gives
\begin{equation}\label{eq:nilpotent-identity}
 (B-1)^2G(r)=\bigl((B-1)^2c+t(B-1)+B\bigr)B^r
                 -(t+r)(B-1)-B.
\end{equation}
Since $B-1$ is a unit modulo $R$ and $R\mid B^r$ for $r\ge E$,
\begin{equation}\label{eq:smooth-endpoint}
 R\mid G(r)\quad\Longleftrightarrow\quad t+r\equiv\rho\pmod R
 \qquad(r\ge E).
\end{equation}
The only omitted positive offsets are $1,\ldots,E-1$, exactly those
tested in $\mathcal E_d$. Because $R\le n<10^d$, we have
$\alpha<4d$ and $\beta<2d$, so $E\le4$.

Now work modulo the odd number $m$. By definition $B^h\equiv1$.
Appending $h$ fixed-width blocks maps the state $(C,y)$ to
$(C+ys+v,y+h)$ modulo $m$. Starting after $j$ blocks and iterating this
affine map $z$ times gives
\[
 G(j+hz)\equiv V_j+z((t+j)s+v)+hs\frac{z(z-1)}2\pmod m.
\]
Multiplication by 2 is invertible modulo $m$, so
\begin{equation}\label{eq:quadratic-equivalence}
 m\mid G(j+hz)\quad\Longleftrightarrow\quad Q_j(z)\equiv0\pmod m.
\end{equation}
Every nonnegative offset has a unique expression $j+hz$ with
$0\le j<h$ and $z\ge0$. Equation~\eqref{eq:smooth-endpoint} becomes
$hz\equiv\rho-t-j\pmod R$, which is soluble exactly when the
stated divisibility by $g$ holds; in that case it is $z\equiv b_j\pmod u$.
As $q\mid m$ and $u\mid R$, their moduli are coprime.
Equation~\eqref{eq:crt-candidate} therefore gives exactly all compatible
solutions $z\equiv w\pmod{qu}$, hence endpoints $x\equiv x_0\pmod P$.

The ceiling in \eqref{eq:first-candidate} selects the least such endpoint
at least $t+E$. If it exceeds $10^d-1$, that progression has no endpoint
in the actual digit interval. Otherwise it is the smallest endpoint of
its progression in the interval. Any selected endpoint has $x-t\ge1$;
since $x-t\equiv j\pmod h$ and $j<h$, its corresponding $z$ is nonnegative.
Thus the use of the iterated-map identity is legitimate.

These implications prove both that every retained candidate is divisible
by $n$ and that every divisible endpoint in this digit interval is either
in $\mathcal E_d$ or no smaller than a candidate in $\mathcal C_d$.
Corollary~\ref{cor:uniform} ensures at least one such endpoint before
$10^n$. Taking the minimum over the stated widths proves
\eqref{eq:minimum-formula}, including exclusion of \emph{all} shorter
lengths, not only those of the same digit width.

\subsection{Explicit quadratic classes, including singular roots}
For completeness, $\mathcal R_j$ need not be treated as an unspecified
congruence-solving operation. The following gives it by modular square
roots and CRT for any quadratic $Az^2+Bz+C$ modulo an odd $m$.
Factor $m=\prod p^e$ and handle each prime power separately.
First divide $A,B,C$ and $p^e$ by their largest common power of $p$.
The resulting congruence is equivalent as a condition on the integer
$z$; its reduced modulus suffices to describe the original solution set.
If the reduced modulus is 1, the sole class is $(0,1)$.
Otherwise write it as $p^f$, so at least one coefficient is a unit.
\begin{enumerate}
\item If $p\mid A,B$, then $p\nmid C$ and there are no roots.
\item If $p\mid A$ and $p\nmid B$, let $y$ be the unique root of
$y^2\equiv B^2-4AC\pmod{p^f}$ satisfying $y\equiv B\pmod p$.
The sole class is
\[
 z\equiv-2C(B+y)^{-1}\pmod{p^f}.
\]
The indicated square root exists and is unique by simple-root lifting.
The denominator is a unit. Substitution verifies the formula;
conversely a root $z$ gives $y=2Az+B$ and
$(B+y)z=-2C$, proving completeness.
\item If $p\nmid A$, put $\delta=[B^2-4AC]_{p^f}$.
For $\delta=0$, there is one class
\[
 z\equiv-B(2A)^{-1}\pmod{p^{\lceil f/2\rceil}}.
\]
For $\delta\ne0$ of odd $p$-adic valuation there are no roots.
Otherwise put $v_p(\delta)=2s<f$. If the unit
$\delta/p^{2s}$ is not a square modulo $p$, there are no roots.
If it is, take either unit square root $y$ modulo $p^{f-2s}$.
The two classes are
\[
 z\equiv(-B\mathbin{\pm}p^s y)(2A)^{-1}\pmod{p^{f-s}}.
\]
These cases follow from the equivalent square congruence
$(2Az+B)^2\equiv\delta\pmod{p^f}$.
\end{enumerate}
Every prime contributes at most two classes, with moduli powers of that
prime. Combining them by CRT gives disjoint classes $(r,q)$ with $q\mid m$
and at most $2^{\omega(m)}$ possibilities. This includes all free higher
digits of singular roots without enumerating them.

\subsection{An example requiring several digit intervals}
For $n=451=11\cdot41$, the formula gives the following complete stable
classes of endpoints. Here $R=1$, so there are no transient exceptions.
\[
\begin{array}{c|c|l}
d&c_d\bmod451&\text{endpoint residue classes}\\\hline
3&0&451,1902,1133,3794,4345,1066,1287,1748,3179,450\pmod{4510}\\
4&241&\varnothing\\
5&196&\varnothing\\
6&63&2084,555,1501,137,478\pmod{2255}
\end{array}
\]
None of the width-3 classes meets $[452,999]$. The width-4 and width-5
quadratics have no compatible roots at all. At width 6, the ceiling
formula gives first endpoints $101304,102030,100721,101612,101953$,
respectively. Hence the least endpoint is $100721$ and
$a(451)=100721-451=100270$. The formula explains the long gap by
whole empty digit intervals, without testing their individual lengths.

\section{First exact family: a last-block principle}
\begin{theorem}[Decimal-smooth moduli]\label{thm:smooth}
For any $n\ge1$ with $D=\ell(n+1)$, if $n\mid10^D$, then
\[\boxed{a(n)=n}.\]
In particular $a(10^r)=10^r$ for all integers $r\ge0$.
\end{theorem}
\begin{proof}
Every $n+k$ for $k\ge1$ has at least $D$ digits, so $n\mid10^{\ell(n+k)}$. The recurrence thus yields $C_n(k)\equiv k\pmod n$ \emph{for all} $k\ge1$, including after any digit-width transitions. All $1\le k<n$ are nonzero modulo $n$, while $k=n$ is zero.
\end{proof}
\begin{corollary}
The formula holds whenever $n=2^u5^v$ and $\max(u,v)\le\ell(n+1)$. Examples are $a(20)=20$, $a(25)=25$, $a(125)=125$, and $a(1000)=1000$.
\end{corollary}
\section{Second exact family: an alternating residue}
\begin{theorem}[An infinite family with $a(n)=2n-1$]\label{thm:seven}
For every integer $j\ge0$ let $D=6j+3$ and $n=(10^D+1)/7$. Then $n$ is a positive integer and
\[
 \boxed{a\left(\frac{10^{6j+3}+1}{7}\right)=\frac{2(10^{6j+3}+1)}7-1}.
\]
In particular, $a(143)=285$.
\end{theorem}
\begin{proof}
Since $10\equiv3\pmod7$, $3^3\equiv-1\pmod7$, and $3$ has order $6$ modulo $7$, we have $7\mid10^D+1$. For $D\ge3$ and $n=(10^D+1)/7$, the inequalities
\[
 10^{D-1}\le n+1\quad\text{and}\quad 3n-1<10^D
\]
hold. Hence all terms $n+1,\ldots,n+(2n-1)$ have exactly $D$ digits. Moreover $10^D=7n-1\equiv-1\pmod n$. Proposition~\ref{prop:width} now yields $C_n(k)\equiv\lceil k/2\rceil\pmod n$ for every $1\le k\le2n-1$. The first value in this range whose ceiling is divisible by $n$ is precisely $k=2n-1$. Thus no shorter prefix is divisible by $n$, and that prefix is.
\end{proof}
\section{A companion exact criterion}
\begin{proposition}[Odd prime powers]\label{prop:primepower}
Suppose $n>1$ is an odd prime power, $D=\ell(n+1)$, $n\mid10^D-1$, and $10^D\ge2n$. Then
\[\boxed{a(n)=n-1}.\]
For example $a(9)=8$ and $a(11)=10$.
\end{proposition}
\begin{proof}
The hypotheses ensure $L=10^D-n-1\ge n-1$. The constant-width recurrence gives $C_n(k)\equiv k(k+1)/2\pmod n$ for $k\le n-1$. Since $n=p^e$ with $p$ odd, $n$ divides $k(k+1)/2$ if and only if $p^e$ divides $k$ or $k+1$, because these consecutive factors are coprime. The first such positive $k$ is $n-1$.
\end{proof}
\section{Prior art, verification, and limitations}
The classification of \emph{all} indices with $a(n)=2$ follows from the two-block congruence $C_n(2)\equiv10^{\ell(n+2)}+2\pmod n$. The equivalent complete formula was published as a formula for OEIS \href{https://oeis.org/A069860}{A069860} by David Wasserman in 2005. We do not present it as a new theorem or claim credit for that classification.

The dependency-free script \texttt{minimum\_formula.py} evaluates
Theorem~\ref{thm:minimum-formula} and exposes the residue classes with
\texttt{--explain}. Its output agrees with all 10,000 terms in the OEIS
b-file. Its quadratic solver is checked exhaustively modulo every odd
integer at most 31, with additional singular and higher-prime-power
controls. Arbitrary-state checks test both directions of the full
residue-class classification. The separate \texttt{least.py} state-orbit
method and direct enumeration provide independent minimum checks.
The script \texttt{verify\_existence.py} independently recomputes
constructive witnesses, and \texttt{verify.py} checks the older families.
Finite checks supplement the proofs. The all-index result is the finite
arithmetic formula \eqref{eq:minimum-formula}; a simpler elementary
expression without finite minimization is not established. Broader
literature priority and external peer-review acceptance remain unverified.
\begin{thebibliography}{9}
\bibitem{prefix} Q. Wang, \emph{Divisibility of Concatenated Arithmetic Progressions and Champernowne Prefixes}, September 11, 2026, repository manuscript, \url{https://github.com/dongxuelian2/oeis-papers/tree/main/papers/A075000-A075002}.
\bibitem{A069862} OEIS Foundation Inc., \emph{A069862: Smallest $k$ such that $n$ divides the concatenation of numbers from $(n+1)$ to $(n+k)$}, \url{https://oeis.org/A069862} (accessed October 22, 2026).
\bibitem{A069860} D. Wasserman, formula for \emph{A069860: Numbers $m$ that divide the concatenation of $m+1$ and $m+2$}, OEIS entry, August 25, 2005, \url{https://oeis.org/A069860}.
\end{thebibliography}
\end{document}
